\chapter[La Mirada del Desprecio]{La Mirada del Desprecio\\ \large Al despreciar el valor ajeno, pierdes tu salud espiritual.}

\elegantimage{11_mirada_desprecio/web/assets/hero.png}

\lettrine[lines=3, nindent=0.5em, findent=0.2em]{M}{}irar a los vulnerables como si fueran basura es un acto de soberbia que envenena la propia salud física y mental. Es ignorar que todos formamos parte de la misma red de vida.

\section{El Espejo Roto}
\begin{elegantquote}
\textit{Cierta persona caminaba siempre mirando por encima del hombro, con la nariz en alto, tratando a los menos afortunados como si fueran basuras molestas. Creía en su propia perfección, incapaz de notar el dolor en los ojos de quienes humillaba...}

\textit{Pero en el inmenso tapiz de Dios, todos somos hilos de la misma tela. Lastimar la existencia del otro es despreciarnos a nosotros mismos. Y así, el universo le llevó a una cama fría, atado a una enfermedad que lo desfiguraba y alejaba a todos a su alrededor.}

\textit{Se convirtió en el ser rechazado que él mismo despreciaba. Sus gritos por una simple caricia humana le rasgó por fin el velo de la soberbia. Comprendió el profundo valor de sentirse igual y hermanos. Al sollozar en su dolorosa soledad, nació la verdadera humildad... que lo haría invencible en el amor verdadero.}

\end{elegantquote}

\vspace{1em}
\begin{center}
    \textbf{\Large La compasión es la medicina del alma.}
\end{center}
\vspace{1em}

\section{Análisis Kármico y Traducción}
\begin{wrapfigure}{r}{0.5\textwidth}
  \vspace{-1.5em}
  \begin{center}
  \inlineelegantimage[0.45\textwidth]{11_mirada_desprecio/web/assets/pasaje_original.jpg}
  \end{center}
  \vspace{-1.5em}
\end{wrapfigure}
La mirada de desprecio y la burla hacia el infortunio ajeno corrompen nuestro ser interior, manifestándose tarde o temprano en enfermedades severas y repulsión social prolongada.

